\documentclass[11pt,a4paper]{article}

\usepackage[utf8]{inputenc}
\usepackage[T1]{fontenc}
\usepackage[brazil]{babel}
\usepackage[margin=2.5cm]{geometry}
\usepackage{graphicx}
\usepackage{booktabs}
\usepackage{amsmath,amssymb}
\usepackage{enumitem}
\usepackage[hidelinks]{hyperref}

\setlength{\parskip}{0.4em}

\title{\textbf{Interpretabilidade de detectores de deepfake de áudio\\
em português brasileiro: plano de pesquisa e metodologia}}
\author{Equipe de pesquisa}
\date{\today}

\begin{document}
\maketitle

\begin{abstract}
\noindent
Este documento descreve o plano do nosso estudo de interpretabilidade de
detectores de \emph{deepfake} de áudio aplicados ao português brasileiro. A ideia
é explicar, de forma simples, o que pretendemos investigar e como pretendemos
fazer isso, do começo ao fim. Usamos uma convenção de registro para manter a
linguagem honesta: chamamos de descritivo o que apenas descreve um padrão nos
dados, de inferencial o que passa por um teste estatístico formal, e de causal
apenas o que se baseia numa intervenção direta no áudio.
\end{abstract}

\section{A pergunta}

Detectores de \emph{deepfake} de áudio de ponta costumam ser treinados em inglês.
A pergunta que orienta este trabalho é dupla e simples: quais pistas acústicas um
detector usa para decidir se um áudio é real ou falso, e como essas pistas mudam
quando o modelo passa por uma adaptação leve ao português brasileiro?

Reduzimos isso a duas perguntas-guia:

\begin{enumerate}[label=\textbf{Q\arabic*.}]
  \item As pistas mudam entre um modelo \emph{zero-shot} e um modelo adaptado?
  \item Se mudam, o que muda (para onde o foco do modelo se desloca no espectro, e
        se isso tem relação com os erros que ele comete)?
\end{enumerate}

Nosso interesse é a interpretabilidade populacional. Não queremos explicar um
clipe isolado, mas descrever, de forma estatística e verificável, o comportamento
do detector sobre um conjunto representativo de áudios.

\section{Visão geral do plano}

Vamos trabalhar com dois detectores construídos sobre o mesmo codificador de fala
congelado, o que isola o efeito da adaptação:

\begin{itemize}
  \item $D_{\mathrm{zs}}$ (\emph{zero-shot}), com a cabeça de decisão original,
        sem qualquer ajuste ao português;
  \item $D_{\mathrm{ad}}$ (adaptado), com uma cabeça leve treinada apenas sobre os
        \emph{embeddings} em português.
\end{itemize}

Sobre esse par, aplicaremos duas análises complementares, ambas na mesma partição
de frequência (8 bandas mel):

\begin{itemize}
  \item uma análise associativa, que mede quais bandas acompanham o score do
        detector;
  \item uma análise causal, que mede quais bandas, ao serem removidas do áudio, de
        fato alteram o score.
\end{itemize}

Por fim, mediremos a convergência entre as duas análises e faremos uma etapa
confirmatória, ligando as pistas mais relevantes à estrutura de erro do detector.
Todo o processo roda como um pipeline em etapas encadeadas:

\begin{quote}
\texttt{collect} $\rightarrow$ \texttt{embeddings} $\rightarrow$ \texttt{adapt}
$\rightarrow$ \texttt{features} $\rightarrow$ \texttt{association} $\rightarrow$
\texttt{occlusion} $\rightarrow$ \texttt{confirmatory} $\rightarrow$ \texttt{report}.
\end{quote}

\section{Dados e amostragem}

Usaremos o corpus BRSpeech-DF (\texttt{AKCIT-Deepfake/BRSpeech-DF}), que reúne
cerca de 459 mil clipes, sendo aproximadamente 77 mil reais (\emph{bonafide}) e
383 mil falsos (\emph{spoof}). O corpus é grande e bastante desbalanceado (cerca
de um real para cada cinco falsos).

Como o objetivo é explicabilidade populacional, e não um \emph{ranking} de
desempenho, montaremos um conjunto de análise equilibrado:

\begin{itemize}
  \item para treinar a cabeça $D_{\mathrm{ad}}$, um subconjunto balanceado de
        1.500 clipes por classe;
  \item para a análise, todo o \emph{bonafide} disponível no teste (1.320 clipes,
        a classe escassa) somado a 1.500 \emph{spoof}, totalizando 2.820 clipes.
\end{itemize}

O conjunto de análise nunca é usado no treino da cabeça, evitando vazamento. Esse
tamanho já oferece poder estatístico para as etapas seguintes, e usar todo o
\emph{bonafide} escasso mantém os quadrantes de decisão significativos. Escalar
para o corpus inteiro fica como trabalho futuro.

\section{Os dois detectores}

O codificador é um modelo wav2vec~2.0 / XLS-R congelado (\emph{checkpoint}
\texttt{nii-yamagishilab/\allowbreak mms-300m-\allowbreak anti-deepfake}), que
produz \emph{embeddings} de 1024 dimensões por clipe. Sobre ele:

\begin{itemize}
  \item $D_{\mathrm{zs}}$ usa a cabeça anti-\emph{spoofing} original do
        \emph{checkpoint};
  \item $D_{\mathrm{ad}}$ acrescenta uma padronização seguida de uma regressão
        logística treinada sobre os \emph{embeddings}, devolvendo uma
        probabilidade de \emph{spoof} calibrada.
\end{itemize}

Congelar o codificador e treinar apenas a cabeça é uma decisão importante, porque
isola o efeito da adaptação no uso das representações. Assim, qualquer mudança que
observarmos refletirá uma re-ponderação de representações que o modelo já possui,
e não a aprendizagem de \emph{features} novas para o português. Por isso, as
conclusões valerão para adaptação rasa, sem extrapolação para um \emph{fine-tuning}
completo.

Antes de qualquer análise de pistas, verificaremos uma pré-condição: a adaptação
precisa de fato melhorar a detecção. Só compararemos as pistas se o modelo
adaptado tiver desempenho melhor que o \emph{zero-shot}. Caso contrário, o
contraste seria reportado como adaptação insuficiente.

\section{Pré-processamento}

Cada clipe será convertido para mono, reamostrado para 16~kHz, recortado em uma
janela fixa de aproximadamente 4~segundos e normalizado. Um ponto importante do
plano: as \emph{features} acústicas serão extraídas exatamente do sinal que entra
no detector, e não do áudio bruto. Dessa forma, a descrição acústica corresponde
ao que o modelo de fato processa, e a normalização evita atribuir pistas ao nível
absoluto do sinal.

\section{A descrição acústica: energia log-mel por banda}

Vamos descrever cada áudio pela energia log-mel em 8 bandas de frequência,
resumindo cada banda pela média e pelo desvio ao longo do tempo, o que dá 16
valores por clipe. Na prática, calculamos o espectro de potência em janelas
curtas, somamos a energia dentro de cada banda, aplicamos o logaritmo e tiramos a
média e o desvio ao longo dos quadros.

A escolha dessa descrição tem um motivo central: ela é localizada em frequência.
Assim, a análise associativa e a análise causal passam a viver no mesmo eixo de
frequência e podem ser comparadas banda a banda, algo que descrições cepstrais
(como MFCC) não permitem de forma direta, por não serem localizadas. As 8 bandas
(em Hz) que usaremos são:

\begin{center}
\small
\begin{tabular}{cccccccc}
\toprule
1 & 2 & 3 & 4 & 5 & 6 & 7 & 8 \\
\midrule
20--282 & 282--639 & 639--1125 & 1125--1788 & 1788--2693 & 2693--3926 & 3926--5607 & 5607--7900 \\
\bottomrule
\end{tabular}
\end{center}

Vale a ressalva de que energia de banda não é fonema. Ela indica onde está a
potência no espectro, não é uma prova fonética.

\section{Como vamos avaliar o desempenho}

A métrica principal de anti-\emph{spoofing} será a EER (\emph{Equal Error Rate}),
o ponto de operação em que a taxa de rejeitar áudio real se iguala à de deixar
passar \emph{deepfake}. Ela resume o detector em um único número, independente do
limiar (quanto menor, melhor), e define o limiar de operação que usaremos no resto
da análise. Como métricas complementares, acompanharemos o coeficiente de
Matthews (MCC) e a acurácia.

Fixando o limiar no ponto de EER e cruzando as predições com os rótulos reais,
formam-se os quatro quadrantes de decisão: \emph{spoof} corretamente detectado,
\emph{bonafide} corretamente aceito, \emph{bonafide} marcado como \emph{spoof} por
engano, e \emph{spoof} que passou despercebido. Esses quadrantes serão a base da
etapa confirmatória.

\section{As hipóteses}

As duas perguntas-guia se desdobram em três hipóteses testáveis, cada uma com um
registro explícito:

\begin{description}
  \item[H1, deslocamento associativo (descritivo).] O perfil de associação entre
        as \emph{features} de banda e o score difere entre os dois detectores.
        Compararemos os perfis de forma descritiva, sem afirmar significância da
        diferença.
  \item[H2, deslocamento causal (causal, com incerteza).] As bandas cuja oclusão
        reduz a probabilidade de \emph{spoof} diferem entre os detectores. Como é
        uma intervenção, reportaremos os perfis com incerteza quantificada
        (intervalos de confiança por \emph{bootstrap}).
  \item[H3, diferenças ligadas ao erro (inferencial).] Para as \emph{features}
        mais relevantes, as distribuições diferem entre os quadrantes de decisão.
        Esta é a única hipótese testada de forma inferencial, com controle de
        falsas descobertas.
\end{description}

\section{Análise associativa (H1)}

Para cada detector, calcularemos a correlação de Spearman (de posto, monotônica)
entre cada \emph{feature} de banda e o score, dentro de cada classe
(\emph{bonafide} e \emph{spoof} separados), com correção de FDR por detector.

Fazer isso dentro de cada classe é essencial. Como o score já separa as classes,
juntar tudo criaria correlação espúria para quase qualquer \emph{feature} que
apenas diferisse entre classes. Dentro da classe, a correlação reflete covariação
genuína. Escolhemos a correlação de Spearman por ser transparente e independente
de modelo, sem a necessidade de um modelo substituto. O custo é que a medida é
univariada, ou seja, não modela interações entre bandas.

\section{Análise causal (H2)}

Aqui vamos intervir no áudio. Removeremos uma banda de frequência por vez, com um
filtro \emph{band-stop} de fase zero, e mediremos a queda na probabilidade de
\emph{spoof}. A leitura é direta: se ocluir a banda derruba o score, aquela banda
contribuía para a decisão de \emph{spoof}; se ocluir a banda aumenta o score, ela
contribuía para a decisão de \emph{bonafide}.

Faremos isso sobre um subconjunto estratificado por quadrante e quantificaremos a
incerteza com intervalos de confiança de 95\% por \emph{bootstrap}. Uma banda cujo
intervalo não inclui o zero indica efeito causal confiável.

Já registramos uma ressalva importante deste método, a ser tratada numa fase
posterior: o filtro remove o conteúdo da banda, mas também introduz uma pequena
distorção. Assim, uma queda pode vir da ausência da banda ou do artefato de
filtragem. Para isolar o efeito do conteúdo, uma etapa futura substituirá a banda
por ruído com energia equivalente.

\section{Convergência das duas análises}

Como as duas análises vivem na mesma grade de bandas, poderemos medir, para cada
detector, o quanto elas concordam: a correlação, banda a banda, entre o perfil
causal (queda de oclusão) e o perfil associativo. Um valor alto indicaria que as
bandas que o detector usa de fato (causal) são também as mais associadas ao score.

Aqui cabe uma ressalva de honestidade: como há apenas 8 bandas, esse número será
um resumo descritivo, e não um teste com poder estatístico. Vamos tratá-lo como
uma indicação, não como prova.

\section{Testes confirmatórios (H3)}

Por fim, aplicaremos testes estatísticos clássicos apenas nas \emph{features} mais
relevantes (para evitar testar as 16 sem critério), ligados diretamente à
estrutura de erro:

\begin{itemize}
  \item um teste de diferença de média (Welch) entre os \emph{bonafide} aceitos e
        os \emph{bonafide} rejeitados por engano, para entender por que o detector
        confunde real com falso;
  \item um teste de diferença de variância (Levene) entre os \emph{spoof} pegos e
        os \emph{spoof} que passaram, para entender por que ele perde \emph{spoofs}.
\end{itemize}

Aplicaremos correção de FDR sobre o conjunto de testes de cada detector. Esta é a
única parte inferencial do estudo, em que caberá a palavra significativo.

\section{O que o estudo vai produzir}

Ao final, esperamos entregar:

\begin{itemize}
  \item uma comparação de desempenho entre os dois detectores (EER, MCC e
        acurácia), que também serve de pré-condição para a análise de pistas;
  \item um perfil associativo por banda para cada detector;
  \item um perfil causal por banda, com intervalos de confiança;
  \item uma medida de convergência entre as duas análises;
  \item um conjunto de testes confirmatórios ligados aos erros do detector.
\end{itemize}

Juntas, essas saídas devem responder às duas perguntas-guia de forma verificável,
com a evidência causal (a oclusão) como parte mais forte, por ser a única baseada
em intervenção.

\section{Ressalvas previstas e trabalho futuro}

\begin{itemize}
  \item Tratar a distorção do filtro de oclusão, substituindo a banda por ruído
        com energia equivalente.
  \item Escalar do subconjunto para o corpus completo e estratificar por gerador
        de \emph{spoof}.
  \item Lembrar que os achados valerão para adaptação rasa, sem extrapolar para um
        \emph{fine-tuning} completo do codificador.
  \item Lembrar que energia de banda indica onde está a potência, não é prova
        fonética.
\end{itemize}

\section{Reprodutibilidade e engenharia}

O estudo roda como um pipeline em etapas resumível. Cada etapa grava seus
artefatos e um marcador de conclusão com o \emph{hash} da configuração, de modo
que reexecutar pula o que já está feito. As sementes aleatórias são fixas e o
ambiente é encapsulado em Docker, com execução automática em GPU. Os parâmetros
planejados para a rodada estão na Tabela~\ref{tab:config}.

\begin{table}[t]
  \centering
  \caption{Parâmetros planejados da rodada.}
  \label{tab:config}
  \begin{tabular}{ll}
    \toprule
    Parâmetro & Valor \\
    \midrule
    Clipes de treino por classe & 1.500 \\
    Clipes de teste por classe & 1.500 (\emph{bonafide} limitado a 1.320) \\
    Bandas de frequência & 8 (de 20~Hz a 7900~Hz) \\
    Clipes por quadrante na oclusão & 150 (cerca de 600 por detector) \\
    Reamostras do \emph{bootstrap} & 1000 \\
    \emph{Features} destacadas para H3 & 10 \\
    Semente aleatória & 42 \\
    \bottomrule
  \end{tabular}
\end{table}

\end{document}
